\section{Referencias}
\label{sec:referencias}

\begin{enumerate}
  \renewcommand{\labelenumi}{[\arabic{enumi}]}

  \item Kit Pagos Colombia, ``Repositorio del proyecto,'' GitHub. [En línea]. Disponible en: \url{https://github.com/PurosBrothers/Kit-Pagos-Colombia---Tesis}

  \item Kit Pagos Colombia, ``Especificación de Requisitos de Software (SRS), versión 2.0,'' \path{docs/final-doc/srs/}, 2026.

  \item Kit Pagos Colombia, ``Documento de Arquitectura de Software (SAD),'' \path{docs/final-doc/sad/}, 2026.

  \item Kit Pagos Colombia, ``API de Simulación: el segundo componente,'' \path{docs/02-arquitectura/3-api-de-simulacion.md}, 2026.

  \item Kit Pagos Colombia, ``Documentación de datos de prueba,'' \path{docs/testing-data/}, 2026.

  \item Kit Pagos Colombia, ``Prototipos del experimento de la Fase 5,'' \path{prototypes/README.md}, 2026.

  \item Kit Pagos Colombia, ``Colección de Postman de la API de Simulación,'' \path{simulator-api/postman/}, 2026.

  \item Meta Open Source, ``Jest: JavaScript testing framework.'' [En línea]. Disponible en: \url{https://jestjs.io/docs/getting-started}

  \item Fastify, ``Fastify documentation.'' [En línea]. Disponible en: \url{https://fastify.dev/docs/latest/}

  \item Postman, ``Postman learning center.'' [En línea]. Disponible en: \url{https://learning.postman.com/docs/}

  \item OpenAPI Initiative, ``OpenAPI Specification v3.0.3,'' 2021. [En línea]. Disponible en: \url{https://spec.openapis.org/oas/v3.0.3}

  \item GitHub, ``GitHub Actions documentation.'' [En línea]. Disponible en: \url{https://docs.github.com/en/actions}

  \item OpenJS Foundation, ``Node.js releases.'' [En línea]. Disponible en: \url{https://nodejs.org/en/about/previous-releases}

  \item Wompi, ``Documentación para desarrolladores.'' [En línea]. Disponible en: \url{https://docs.wompi.co/}

  \item Mercado Pago, ``Documentación para desarrolladores.'' [En línea]. Disponible en: \url{https://www.mercadopago.com.co/developers/es/docs}

  \item Kushki, ``Documentación para desarrolladores.'' [En línea]. Disponible en: \url{https://docs.kushkipagos.com/}

  \item Rapyd, ``Rapyd API documentation.'' [En línea]. Disponible en: \url{https://docs.rapyd.net/}

  \item M. Collina \emph{et al.}, ``autocannon: fast HTTP/1.1 benchmarking tool written in Node.js,'' versión 8.0.0. El histograma de latencia se define en el archivo \emph{histUtil.js} de la carpeta \emph{lib}, con un valor mínimo distinguible de 1~ms, y \emph{run.js}, en la misma carpeta, solo registra las latencias de respuestas 2xx. [En línea]. Disponible en: \url{https://github.com/mcollina/autocannon}\label{ref:autocannon}

  \item Fastify, ``Fastify benchmarks.'' [En línea]. Disponible en: \url{https://github.com/fastify/benchmarks}\label{ref:fastify-benchmarks}

  \item OpenJS Foundation, ``Performance measurement APIs (perf\_hooks),'' documentación de Node.js. [En línea]. Disponible en: \url{https://nodejs.org/api/perf_hooks.html}\label{ref:node-perf-hooks}

  \item D. Eastlake y A. Panitz, ``Reserved Top Level DNS Names,'' RFC 2606, 1999. [En línea]. Disponible en: \url{https://www.rfc-editor.org/rfc/rfc2606.txt}\label{ref:rfc2606}

  \item IEEE, ``IEEE Standard for Software and System Test Documentation,'' IEEE Std 829-2008, 2008.

  \item PCI Security Standards Council, ``PCI DSS: Payment Card Industry Data Security Standard.'' [En línea]. Disponible en: \url{https://www.pcisecuritystandards.org/}
\end{enumerate}
